\section{A different certificate mechanism: inert-prime CRT closure}
\label{sec:inert-crt}
The following is a new compiler-interface theorem within this comparison,
not a new measured consistency polynomial. It changes the arithmetic
certificate for bounded universal quantification, rather than the source
Turing machine or a Jones packing weight. Its complete proof is elementary.
The factorial, binomial and variable-exponent graphs have not been expanded
and counted in a final ordinary polynomial; that limitation is essential.

\paragraph{Input.}
Fix $n\geq1$, natural parameters $x$, and polynomials
$e(i,w;x),f(i,w;x)$ with nonnegative integer coefficients, where
$w=(w_1,\ldots,w_m)$ and $m\geq1$. The target is
\[
 \forall i<n\;\exists w\in\mathbb N^m\quad e(i,w;x)=f(i,w;x).
\]
An integer-coefficient polynomial can first be split into positive and
negative parts. For a nonempty fixed list, define a binary-tree norm by
$\mathcal N([z])=z$ and
$\mathcal N(L\mathbin{\|}R)=\mathcal N(L)^2+\mathcal N(R)^2$, with both
children nonempty.

\paragraph{Explicit certificate.}
Existentially choose $U,A_1,\ldots,A_m,c\in\mathbb N$. Set
\begin{align*}
H&=e(n,U,\ldots,U;x)+f(n,U,\ldots,U;x)+n+U+1, & T&=H!,\\
Q&=(4Tn)^n+1, & M&=4TQ+1,\\
B&=\bigl((4T)^n n!\tbinom{Q+n}{n}\bigr)\bmod M,\\
z_j&=\tbinom{A_j}{U+1}\quad(1\leq j\leq m),&
z_{m+1}&=e(c,A;x)-f(c,A;x),\\
z_{m+2}&=4T(c+1)-1.
\end{align*}
Then the target is equivalent to the single condition
\[
 B\mid\mathcal N(z_1,\ldots,z_{m+2}).
\]
Here $\binom{A}{k}=0$ for natural $A<k$. Any future graph replacement
must implement that same total convention.
These are an explicit special-function certificate, \emph{not yet ordinary
polynomial equations}. Their functions can be replaced by supplied
Diophantine graphs, but those graph witnesses must be counted.

\paragraph{Product and prime facts.}
Put $m_i=4T(i+1)-1$ for $0\leq i<n$. Then
\[
 (4T)^n n!\binom{Q+n}{n}
 =\prod_{j=1}^n(4TQ+4Tj)
 \equiv\prod_{j=1}^n(4Tj-1)\pmod M.
\]
The positive product is less than $(4Tn)^n<Q<M$, so
$B=\prod_i m_i$. Distinct moduli are coprime: a common divisor divides
$j-i$, is coprime to $4T$, and divides $T=H!$ because $j-i\leq n\leq H$.
Each $m_i\equiv3\pmod4$ has a prime divisor $p_i\equiv3\pmod4$.
Every prime divisor of every $m_i$ exceeds $H$, since a prime at most $H$
divides $T$.

\paragraph{Soundness.}
Over $\mathbb F_p$ for $p\equiv3\pmod4$, the equation $a^2+b^2=0$
forces $a=b=0$: otherwise it would make $-1$ a square.
Induction through the norm tree makes every leaf zero modulo $p_i$.
The index leaf gives $c\equiv i\pmod{p_i}$.
The binomial leaves, multiplied by the invertible $(U+1)!$, give
\[
 \prod_{r=0}^{U}(A_j-r)\equiv0\pmod{p_i}.
\]
Primality supplies $w_j\in[0,U]$ congruent to $A_j$ modulo the
\emph{same} prime $p_i$. Hence the local residual at $(i,w;x)$ vanishes
modulo $p_i$. Its absolute value is at most
$e(n,U,\ldots,U;x)+f(n,U,\ldots,U;x)<H<p_i$, by monotonicity.
It is therefore exactly zero.

\paragraph{Completeness.}
Choose all finitely many local witnesses, and let $U$ bound them.
The CRT supplies $c\equiv i$ and $A_j\equiv w_{ij}\pmod{m_i}$.
Each $(U+1)!$ is invertible modulo every $m_i$, so the binomial leaves
vanish modulo $m_i$; so does the local residual. Thus their norm and
the index residual vanish modulo $B$. This proves both directions
without assuming consistency or invoking a standard-model soundness claim.
The argument is finitary and formalizable in ZFC.

\paragraph{Why this is different from the incoming constructor.}
The incoming one-stack source has five local parameters, but its proposed
global constructor first introduces 84 nonlinear wires, then applies
bounded closure to 89 local witnesses. Apply bounded closure to the
original local polynomial instead, and delay degree reduction. Independently,
the present theorem replaces separate coordinate bounds by one $U$, cancels
all per-coordinate factorial factors, and combines $m$ range congruences
and the evaluation and index congruences into a single norm congruence.
At the certificate interface this removes $m-1$ bound coordinates and
$m+1$ separate congruence tests. In an implementation using one integer
quotient per congruence this also saves $m+1$ quotient slots, but natural-domain
graphs may have different bookkeeping. It still needs $m$ binomial graphs.
These are not final Diophantine-variable savings; omitted wires can return
when the final equation is flattened.

A balanced norm has degree at most $2^{\lceil\log_2(m+2)\rceil}$ in its
leaves. Alternatively place the high-degree local residual immediately
below the root, so that it is only squared once, and allow more depth for
the range leaves. Shared arithmetic DAGs should be preserved throughout:
naive substitution into expression trees can overwhelm the savings.

\paragraph{Validation and remaining work.}
The authored checker verifies 28 exact product/coprimality cases,
2,030 anisotropy pairs, 111 local range/recovery cases and a CRT completeness
fixture. Three negative controls demonstrate why ordinary sums of three
squares, primes $1\bmod4$, and omitted residual bounds are invalid
substitutes. These tests supplement rather than prove the theorem.
The artifact does not provide a newly kernel-checked closure or a new
ordinary-polynomial record. Full graph instantiation, coefficient output,
final counts and endpoint calibration remain necessary.
